\scanpage{161}
\mspage{77}
% archive sheet number: 160

par loc. cit. \deleted{\unsure{§}}9.18. b)\inserted{ (p. 156)}, on a

\[
\operatorname{Ind}_{\pi_1}(M_{\pi_1})\xrightarrow{\approx}M,
\]

ce qui prouve d).

\uline{Prouvons $\mathrm{d})\Longrightarrow\mathrm{PTA}$} Par loc. cit \deleted{\unsure{§}}9.\unsure{18|8} a)\inserted{ (p. 156)}, \deleted{\unsure{on sait}} on voit que les objets de $M_{\pi_1}$ sont $\pi_1$-accessibles dans $\operatorname{Ind}_{\pi_1}(M_{\pi_1})$, \deleted{\unsure{donc}} d'autre part il est clair que $M_{\pi_1}$ est une sous-catégorie pleine strictement génératrice de \deleted{\unsure{$\mathcal C$|$M$}} $\operatorname{Ind}_{\pi_1}(M_{\pi_1})$, donc il reste à voir que $\operatorname{Ind}_{\pi_1}(M_{\unsure{1|\pi_1}})$ est stable par petites $\varinjlim$, i.e. le

\uline{Lemme} 4.3.\deleted{\unsure{6}}7. Soit $\unsure{M|A}_\pi$ une \inserted{\unsure{$U$|$\mathcal U$}-}catégorie stable par limites inductives de cardinal $\leq\pi$, où $\pi$ est un cardinal \inserted{infini} donné. Alors $\operatorname{Ind}_{\pi}M_{\pi}$ est stable par \uline{petites} limites inductives. De plus, \inserted{$M_\pi\longrightarrow M$} $M=\operatorname{Ind}_{\pi}\unsure{C|M_\pi}$ est solution d'un problème 2-universel, à \unsure{savoir} que \deleted{pour} \deleted{\unsure{$M$}} pour tte \deleted{\unsure{$U$-}}catégorie $N$ stable par petites limites inductives, le foncteur restriction \deleted{\unsure{ind}} induit une équivalence

\[
\underline{\operatorname{Hom}}_{!}(M,N)\xrightarrow{\approx}
\underline{\operatorname{Hom}}_{!\pi}(M_{\pi},N),
\tag*{(4.\unsure{3|8}.7.1)}
\]

où \unsure{le|les} \deleted{\unsure{prem}} \unsure{signe|signes} $!$ (resp. $!\pi$) désigne la sous-catégorie pleine de la catégorie $\underline{\operatorname{Hom}}$, formée des foncteurs qui commutent aux petites $\varinjlim$ (resp. aux limites inductives de cardinal $\leq\pi$).

